\input python_header.tex

% python_header.tex defines the language `Python' with JavaScript's keywords; this guide has its own.
\lstdefinelanguage{PythonStyle}{
  keywords={and, as, assert, break, continue, def, elif, else, for, from, global, if, import,
    in, is, lambda, nonlocal, not, or, pass, return, while},
  keywords=[2]{True, False, None},
  morecomment=[l]{\#},
  morestring=[b]',
  morestring=[b]",
  sensitive=true
}
\lstset{language=PythonStyle, keywordstyle=\bfseries, keywordstyle=[2]\bfseries,
  commentstyle=\itshape}

\parindent=0pt
\parskip=0.6em

\begin{document}
  \thispagestyle{empty}

  \markright{SICP, Python Adaptation, Python Style Guide, 2026}
  \begin{center}
    {\Large {\bf Python Style Guide}---2026 edition}\\[10mm]

    {\large Martin Henz}\\[5mm]

    {\large National University of Singapore \\
    School of Computing}\\[10mm]

    {\large \today}\\[10mm]
  \end{center}

\noindent
This is the style guide for the Python sublanguages (\S 1--4) used in the book
\emph{Structure and Interpretation of Computer Programs, Python Edition} (SICPy).
It follows \href{https://peps.python.org/pep-0008/}{PEP~8}, the style that most Python programmers use,
so what you practice here carries over to your future courses and to shared codebases,
and it adds what SICPy needs.
In the beginning of a first-year course, instructors may just \emph{encourage}
adherence to the style, and then gradually enforce adherence in order to
improve communication between learners and staff.

The point of a style is that \textbf{a reader sees the structure of your program at a glance}.
In Python, the structure is also what the interpreter reads: indentation and line breaks are part
of the syntax. Many of the errors that beginners take for logic errors in their programs are
formatting errors. Follow this guide, and you avoid them.

\section*{Indentation}

The body of a function, and of an \lstinline{if}, \lstinline{elif}, \lstinline{else},
\lstinline{while} or \lstinline{for} statement, is indented by \textbf{four spaces} more
than the line that ends with the colon. Use spaces, never tabs, and never a mixture of the two.
(In the Source Academy editor, the Tab key inserts spaces.)

\begin{lstlisting}
# good style
def make_adder(x):
    def the_adder(y):
        if x >= 0:
            return x + y
        else:
            return -x + y
    return the_adder
\end{lstlisting}

\begin{lstlisting}
# bad style: indentation of two spaces (legal, but hard to read)
def make_adder(x):
  def the_adder(y):
    return x + y
  return the_adder
\end{lstlisting}

\begin{lstlisting}
# illegal Python: inconsistent indentation
def make_adder(x):
    def the_adder(y):
        if x >= 0:
            return x + y
          else:
            return -x + y
    return the_adder
\end{lstlisting}

Python reports an \lstinline{IndentationError} when the indentation of a line does not
match any line above it. The cure is to indent every block by the same amount, as in the good example.

\section*{Line length}

A line should have \textbf{at most 79 characters}. In the playground, the vertical line in the
editor marks the limit: no line should cross it. Long lines are hard to read, and hard to compare
when you work side by side with others.

\subsection*{How to break a line}

Python ends a statement at the end of the line.
\textbf{You can break a line only inside parentheses \lstinline{( )}, brackets \lstinline{[ ]}
or braces \texttt{\{ \}}.}
Anywhere else, the line break ends the statement, and Python reports an error such as
``Unexpected token'' at the end of the line.

\begin{lstlisting}
# illegal Python: the statement ends after the +
total = first_amount +
        second_amount
\end{lstlisting}

\begin{lstlisting}
# good style: the break is inside parentheses
total = (first_amount +
         second_amount)
\end{lstlisting}

So when a long expression must be broken, put it in parentheses. This is why you will see
\lstinline{return (} followed by an expression on several lines: the parentheses are what allows
the line breaks (see \emph{Conditional expressions} below). Do not put parentheses around short
expressions: write \lstinline{return x + 1}, not \lstinline{return (x + 1)}.

Break lines after commas in arguments and parameters, and before \lstinline{if} and
\lstinline{else} in conditional expressions. When you break a long expression before or after a
binary operator, do it the same way throughout your program.

\section*{Spaces}

Put a single space on both sides of a binary operator, and after each comma.

\begin{lstlisting}
# good style
x = 1 + 1
y = (x + 2) * 3
message = "zero" if x == 0 else "not zero"
my_function(1, 2, 3)
\end{lstlisting}

\begin{lstlisting}
# bad style
x=1+1
y=(x+2)*3
message = "zero" if x==0 else "not zero"
my_function(1,2,3)
\end{lstlisting}

Do not put a space between a unary minus and its operand: write \lstinline{-x}, not
\lstinline{- x}. (This differs from the style guide of Source, the JavaScript sublanguage.)
Do not put spaces \emph{inside} parentheses, brackets and braces, or \emph{before} the parentheses
of a function application or a function definition.

\begin{lstlisting}
# good style
def my_function(arg1, arg2, arg3):
    return [arg1, arg2, arg3]

result = my_function(1, 2, 3)
negative_x = -x
\end{lstlisting}

\begin{lstlisting}
# bad style
def my_function (arg1, arg2, arg3):
    return [ arg1, arg2, arg3 ]

result = my_function ( 1, 2, 3 )
negative_x = - x
\end{lstlisting}

\paragraph{Trailing whitespace.} There should be no whitespace (spaces or tabs) before a newline:
you cannot see it, and it causes problems in tools that handle program text, such as version control
systems. Clean up all such \emph{trailing whitespace} before you submit your program.

\section*{Blank lines}

Separate the top-level function definitions of your program by two blank lines. Inside a function,
use a single blank line, sparingly, to separate groups of statements.

\begin{lstlisting}
# good style
def square(x):
    return x * x


def cube(x):
    return x * square(x)


result = cube(3)
\end{lstlisting}

\section*{Function definitions}

In \lstinline{def name(parameters):}, there is no space before the opening parenthesis and none
before the colon. The body starts on the next line, indented by four spaces.

\section*{Conditional statements}

Write \lstinline{if}, \lstinline{elif} and \lstinline{else} each at the start of a line, followed by
an indented block on the next lines. Do not put the body after the colon, even if it consists of
one statement.

\begin{lstlisting}
# good style
def sign(x):
    if x > 0:
        return 1
    elif x < 0:
        return -1
    else:
        return 0
\end{lstlisting}

\begin{lstlisting}
# bad style: body on the same line as the condition
def sign(x):
    if x > 0: return 1
    elif x < 0: return -1
    else: return 0
\end{lstlisting}

Do not put parentheses around the condition: write \lstinline{if x > 0:}, not
\lstinline{if (x > 0):}.

\section*{Conditional expressions}

If a conditional expression fits in the line, keep it there.

\begin{lstlisting}
# good style
def absolute(x):
    return x if x >= 0 else -x

aspect_ratio = 4 / 3 if landscape else 3 / 4
\end{lstlisting}

If it does not fit, put it in parentheses, and put the parts that start with \lstinline{if} or
\lstinline{else} on lines of their own, aligned with the first part. Do not use more lines than
needed.

\begin{lstlisting}
# good style
def ackermann(x, y):
    return (
        0 if y == 0
        else 2 * y if x == 0
        else 2 if y == 1
        else ackermann(x - 1, ackermann(x, y - 1))
    )
\end{lstlisting}

\begin{lstlisting}
# bad style: line too long
def ackermann(x, y):
    return 0 if y == 0 else 2 * y if x == 0 else 2 if y == 1 else ackermann(x - 1, ackermann(x, y - 1))
\end{lstlisting}

\begin{lstlisting}
# bad style: too much indentation
def ackermann(x, y):
    return (
        0 if y == 0
        else 2 * y if x == 0
            else 2 if y == 1
                else ackermann(x - 1, ackermann(x, y - 1))
    )
\end{lstlisting}

\begin{lstlisting}
# illegal Python: the line break is not inside parentheses
def ackermann(x, y):
    return 0 if y == 0
    else 2 * y if x == 0
    else ackermann(x - 1, ackermann(x, y - 1))
\end{lstlisting}

\section*{Arguments and parameters}

The arguments of function applications, and the parameters of function definitions and lambda
expressions, are written in one line, if they fit. If they do not fit, they are aligned vertically:
either under the first argument (after the opening parenthesis), or indented by four spaces on the
lines after the opening parenthesis, which is then the last character of its line.

\begin{lstlisting}
# good style: aligned under the first argument
def execute_application(fun, args, succeed, fail):
    return succeed(apply_primitive_function(fun, args),
                   fail)
\end{lstlisting}

\begin{lstlisting}
# good style: indented by four spaces
def make_execution_function(
    inst, labels, machine, pc, flag, stack, ops
):
    return assemble_instruction(
        instruction_source(inst),
        labels,
        machine,
        pc,
    )
\end{lstlisting}

\begin{lstlisting}
# bad style: the arguments are neither aligned nor indented
result = accumulate(add, 0, map(square,
    enumerate_interval(1,
        10)))
\end{lstlisting}

When the arguments are on lines of their own, the closing parenthesis goes on a line of its own,
aligned with the start of the line that has the opening parenthesis (as above), or at the end of
the last argument.

\section*{Lambda expressions}

Keep the parameters and the body expression of a lambda expression in one line, if they are short
enough. If they are lengthy, put the lambda expression in parentheses, break the line after the
colon, and indent the body by four spaces.

\begin{lstlisting}
# good style
def count_buttons(garment):
    return accumulate(
        lambda sleaves, total: sleaves + total,
        0,
        map(
            lambda jacket: (
                count_buttons(jacket) if is_checkered(jacket)
                else 1
            ),
            garment,
        ),
    )
\end{lstlisting}

\section*{Import directives}

If the imported names do not fit in one line, put them in parentheses, indent them by four
spaces, and put the closing parenthesis on a line of its own.

\begin{lstlisting}
# good style
from rune import (
    beside, heart,
    show
)
show(beside(heart, heart))
\end{lstlisting}

\section*{Names}

When naming constants or variables, choose names that are meaningful to your reader and that make
sense in the application domain. Names contain only lower-case letters of the roman alphabet,
digits, and underscores. Use underscores to separate words (\emph{snake case}). Avoid the
single letters \lstinline{l} (lower-case L), \lstinline{O} and \lstinline{I}, which look like digits.

\begin{lstlisting}
# good names
my_variable
x_1
draw_connected_proportional

# bad names
myVariable
Draw_Connected_Proportional
l
\end{lstlisting}

\section*{Nested scope}

Do not use the same name in a nested scope. Python does not warn you, and you can easily change the
meaning of a name by accident.

\begin{lstlisting}
# bad program
x = 1


def f(x):
    # here, the name x of the global frame is
    # "shadowed" by the parameter x
    return x + 1
\end{lstlisting}

\section*{Comments}

A comment starts with \lstinline{#} and a space. Write a comment on a line of its own, above the
code that it explains and at the same indentation, and say \emph{why}, not \emph{what}.

\begin{lstlisting}
# good style
# The accumulator is the sum of the elements seen so far.
def sum_iter(xs, acc):
    return acc if is_null(xs) else sum_iter(tail(xs), acc + head(xs))
\end{lstlisting}

\section*{If your program does not run}

Many error messages of the Python interpreter come from formatting. Check, in this order:

\medskip
\noindent
\begin{tabular}{p{0.34\textwidth}p{0.58\textwidth}}
\hline
\textbf{Message or symptom} & \textbf{Likely cause}\\
\hline
``Unexpected token'' at the end of a line & A line break outside parentheses: put the expression in parentheses, or keep it on one line.\\
\texttt{IndentationError} & Inconsistent indentation: indent every block by four spaces, and do not mix tabs and spaces.\\
``Unexpected token: indent'' & A line is indented that should not be: line it up with the previous statement.\\
A missing colon & Every \texttt{def}, \texttt{if}, \texttt{elif}, \texttt{else}, \texttt{while} and \texttt{for} line ends with a colon.\\
A block that seems to end too early or too late & The indentation of its lines differs: align them under each other.\\
\hline
\end{tabular}

\medskip
\noindent
A program that follows this guide can still be wrong, of course. But if it is, you can be sure that
the reason is the logic and not the layout.

\section*{How this guide differs from the style guide of Source}

Colleagues who know the style guide of the JavaScript sublanguages of Source will recognize most of
the rules. These are the differences that result from Python:

\begin{itemize}
\item Blocks are not enclosed in braces: the indentation \emph{is} the block.
\item A conditional expression is written \lstinline{consequent if predicate else alternative}.
\item A unary minus is written without a space: \lstinline{-x}.
\item A line can be broken only inside parentheses, brackets and braces; a conditional expression in a
      \lstinline{return} statement is therefore put in parentheses when it does not fit in a line.
\item There are two blank lines between top-level function definitions, as in PEP~8.
\end{itemize}

\end{document}
